% Wiley NJD v5 Template - Disease Early Warning and Surveillance
\documentclass[AMA,LATO1COL]{WileyNJD-v2}

% Additional packages
\usepackage{amsmath}
\usepackage{graphicx}
\usepackage{booktabs}
\usepackage{array}
\usepackage{hyperref}
\usepackage{url}

% Article metadata
\articletype{Research Article}

\received{XX Month XXXX}
\revised{XX Month XXXX}
\accepted{XX Month XXXX}

\raggedbottom

\begin{document}

\title{Blockchain-Based Architectural Model for Disease Early Warning and Surveillance}

\author[1]{Teklhiwot Mekonen\textsuperscript{*}}
\author[1]{Esubalew Alemneh}

\authormark{Mekonen and Alemneh}

\address[1]{\orgdiv{Faculty of Computing, Bahir Dar Institute of Technology}, \orgname{Bahir Dar University}, \orgaddress{\city{Bahir Dar}, \country{Ethiopia}}}

\corres{*Teklhiwot Mekonen, Faculty of Computing, Bahir Dar Institute of Technology, Bahir Dar University, Bahir Dar, Ethiopia. \email{teklhiwot@example.com}}

\presentaddress{Faculty of Computing, Bahir Dar Institute of Technology, Bahir Dar University, Bahir Dar, Ethiopia}

\abstract{Disease early warning and surveillance (DEWS) systems in low- and middle-income countries (LMICs) face persistent architectural failures: data centralization, metadata-level privacy exposure, unscalable ledger replication, absence of selective disclosure, and exclusion of resource-constrained peripheral sites from meaningful system participation. Existing blockchain-health architectures exacerbate these weaknesses by treating the ledger as a bulk data repository, assuming uniform node capacity, and conflating immutability with privacy. This paper presents a hybrid, privacy-layered surveillance architecture that replaces full on-chain storage with selective anchoring: encrypted surveillance payloads reside off-chain via IPFS, while Ethereum smart contracts retain only cryptographic hashes, access-control state, provenance records, alerts, and audit events. The architecture introduces five integrated novelty contributions: (1)~hierarchical event processing with district-level batching aligned to public-health administrative tiers; (2)~a layered privacy model combining zero-knowledge verification, dynamic grant/revoke/expire permissions, and encrypted off-chain records; (3)~edge and fog preprocessing to reduce bandwidth, latency, and compute burden at peripheral facilities; (4)~a comparative consensus framework evaluating Proof of Authority, Delegated Proof of Stake, PBFT variants, and Raft-based alliance models against LMIC workload, trust, and infrastructure constraints; and (5)~FHIR-aware interoperability workflows supporting standards-compliant selective disclosure across organizational boundaries. A functional prototype validates the core smart-contract coordination layer. Performance evaluation demonstrates stable block generation (12--14\,s), usability assessment with 26 domain professionals yields Software Usability Measurement Inventory scores significantly exceeding standardized benchmarks, and an expanded threat model addresses validator collusion, metadata leakage, overbroad tier exposure, and equity-aware governance. The architecture advances the state of the art by treating privacy, consensus, and resource asymmetry as linked governance concerns rather than isolated engineering parameters, yielding a scalable, privacy-preserving, and LMIC-optimized surveillance infrastructure.}

\keywords{blockchain, disease early warning, surveillance, Ethereum, smart contracts, health informatics, distributed systems}

\jnlcitation{\cmark[Authors]}

\maketitle

%% ===== INTRODUCTION =====
\section{Introduction}\label{sec:introduction}

Disease early warning and surveillance (DEWS) systems underpin outbreak detection, reporting, and response across hierarchically organized health institutions---from community health posts and primary facilities through district, zonal, and regional offices to national public health authorities. In low- and middle-income countries (LMICs) such as Ethiopia, these systems confront persistent structural failures: data and process centralization that introduces single-point-of-failure risks, security vulnerabilities arising from the absence of tamper-evident record-keeping, traceability gaps across the reporting hierarchy, reliance on manual completeness verification, and interoperability barriers that fragment information flow across organizational boundaries.

Blockchain technology has been proposed as an architectural remedy, offering decentralization, immutability, cryptographic integrity, and transaction traceability.\cite{yaga2018blockchain} However, first-generation blockchain-health architectures---including prior work on DEWS---exhibit four categories of shortcoming that limit their scientific and practical contribution.

First, \emph{monolithic on-chain storage} requires every participating node to replicate the full surveillance payload, producing unsustainable storage growth, high gas costs, and bandwidth demands incompatible with LMIC infrastructure.\cite{kelly2019key} Second, \emph{narrow privacy models} equate immutability with confidentiality, ignoring metadata-level exposure on public or permissioned ledgers, overbroad access across health-system tiers, and the absence of selective disclosure or permission revocation.\cite{finlayson2019adversarial} Third, \emph{unargued consensus assumptions}---typically Proof of Authority adopted without comparative evaluation---leave validator governance, collusion resilience, and trust concentration unaddressed.\cite{core2012building} Fourth, \emph{resource-uniform participation models} assume that all facilities, including rural health posts with intermittent connectivity and underpowered devices, can operate as full blockchain peers, thereby excluding the very populations surveillance is designed to protect.

These limitations motivate a fundamental architectural shift: from full ledger replication to \emph{selective anchoring}, where encrypted surveillance payloads reside off-chain while the blockchain retains only cryptographic hashes, permissions, provenance records, alerts, and audit events.

This paper presents a hybrid, privacy-layered, hierarchical surveillance architecture designed for LMIC deployment contexts. The architecture reframes Ethereum from a bulk data repository to a coordination and proof layer, integrating the following novelty contributions:

\begin{enumerate}
\item A \emph{selective-anchoring architecture} with hierarchical event processing: community and facility nodes submit encrypted case data to off-chain storage (IPFS), district and regional layers batch and validate submissions, and smart contracts commit only hashes, access-control state, and audit artifacts on-chain---eliminating full ledger replication of surveillance payloads.
\item A \emph{layered privacy model} combining zero-knowledge verification for privacy-preserving validation, dynamic smart-contract-based access control with grant, revoke, and expire semantics, and encrypted off-chain records---replacing immutability-only privacy assumptions.
\item \emph{Edge and fog preprocessing} at peripheral tiers to reduce bandwidth, latency, and compute burden, enabling resource-constrained facilities to participate without operating as heavy validator nodes.
\item A \emph{comparative consensus framework} evaluating Proof of Authority, Delegated Proof of Stake, PBFT variants, and Raft-based alliance models against LMIC workload characteristics, trust topology, and infrastructure constraints---with explicit validator governance, appointment, and failure-handling specifications.
\item \emph{FHIR-aware interoperability workflows} supporting standards-compliant selective disclosure and cross-organizational data exchange within the privacy architecture.
\item A functional prototype validating the smart-contract coordination layer, evaluated across three dimensions: performance (block generation timing), usability (Software Usability Measurement Inventory with 26 domain professionals), and security (expanded threat model addressing validator collusion, metadata leakage, overbroad tier exposure, and equity-aware governance).
\end{enumerate}

The remainder of this paper is organized as follows. Section~\ref{sec:background} presents background context and related work. Section~\ref{sec:methodology} describes the research methodology. Section~\ref{sec:model} details the proposed architecture. Section~\ref{sec:implementation} presents the prototype implementation. Section~\ref{sec:evaluation} reports evaluation results. Section~\ref{sec:discussion} discusses implications and limitations. Section~\ref{sec:conclusion} concludes the paper and identifies future research directions.

%% ===== BACKGROUND AND RELATED WORK =====
\section{Background and Related Work}\label{sec:background}

\subsection{Blockchain Foundations for Health Surveillance}\label{subsec:blockchain}

Blockchain provides four properties directly relevant to distributed health surveillance: immutability of committed records, decentralized validation without trusted intermediaries, consensus-driven Byzantine fault tolerance, and end-to-end transaction traceability.\cite{yaga2018blockchain} However, these properties alone are insufficient for privacy-sensitive, resource-constrained surveillance deployments. Full ledger replication forces every node to store the complete surveillance payload, producing unsustainable storage growth and bandwidth demands in low- and middle-income settings. Immutability, while protecting against retroactive tampering, does not prevent metadata exposure, overbroad read access across health-system tiers, or the absence of selective disclosure. Consequently, modern blockchain-health architectures have shifted from monolithic on-chain designs toward \emph{hybrid} models that confine the ledger to coordination artifacts---hashes, permissions, provenance, alerts, and audit events---while encrypted bulk data resides off-chain.

\subsection{Platform and Consensus Rationale}\label{subsec:platform}

Ethereum was selected as the coordination layer based on its native smart-contract support, Turing-complete programmability (Solidity), and mature tooling for permissioned deployment (Table~\ref{tab:platform}). Critically, the architecture reframes Ethereum from a bulk data repository to a \emph{selective-anchoring} layer: smart contracts manage only cryptographic hashes, access-control state, and audit artifacts, while encrypted surveillance payloads are stored off-chain via IPFS. This separation addresses the scalability, cost, and privacy limitations inherent in direct Layer-1 data recording.

\begin{table}[!ht]
\caption{Blockchain platform comparison (selection criteria).}\label{tab:platform}
\begin{tabular*}{\columnwidth}{@{\extracolsep\fill}lccc@{}}
\toprule
Criterion & Ethereum & Bitcoin & XRPL \\
\midrule
Smart contracts & Yes & No & No \\
dApp support & Yes & No & No \\
Turing-complete language & Yes & No & No \\
Digital identity & Yes & Yes & No \\
Data oracles & Yes & Yes & No \\
Human-readable addresses & Yes & Yes & No \\
Block time & 14 s & 10 min & 4 s \\
\bottomrule
\end{tabular*}
\end{table}

Consensus selection is treated as a workload-specific governance decision rather than an assumed default. While the prototype employs Proof of Authority (PoA) for its low latency and deterministic finality, Section~\ref{sec:model} presents a comparative evaluation of PoA against Delegated Proof of Stake (DPoS), Practical Byzantine Fault Tolerance (PBFT) variants, and Raft-based alliance models, assessed against LMIC workload characteristics, trust topology, validator appointment, collusion resilience, and infrastructure constraints.

\subsection{Disease Early Warning and Surveillance in Ethiopia}\label{subsec:dews}

Ethiopia's DEWS system operates through a five-tier hierarchy: community health posts, woreda health offices, zone health offices, regional health offices, and the national Ethiopian Public Health Institute (EPHI)/Public Health Emergency Management (PHEM) center. The system conducts indicator-based surveillance (structured scheduled reporting) and event-based surveillance (rapid informal signal capture). Twenty priority diseases are classified into 13 immediately reportable conditions (notification within 30 minutes) and 7 weekly reportable conditions, with report completeness calculated weekly at every tier. Laboratory surveillance operates through a parallel tiered network (woreda, zone, regional, national) responsible for specimen confirmation.

This administrative hierarchy maps directly to the proposed multi-tier computing architecture. Peripheral health posts and community workers operate in environments characterized by intermittent connectivity, underpowered devices, and limited technical capacity---conditions that preclude full blockchain peer responsibilities. District and zonal offices possess intermediate infrastructure suitable for batching, validation, and edge/fog preprocessing. Regional and national tiers command sufficient resources for validator duties and analytical workloads. The architecture therefore assigns asymmetric node roles aligned to each tier's actual capacity, ensuring that resource-constrained peripheral sites participate as lightweight reporting nodes rather than heavy consensus participants.

\subsection{Related Work}\label{subsec:related}

\subsubsection{Blockchain-Based Healthcare Architectures}

Blockchain adoption in healthcare has progressed through identifiable architectural generations. Early systems stored clinical records directly on-chain: Omar et al.\cite{omar2019privacy} proposed a patient-centric blockchain using decentralization to mitigate cloud security risks, but lacked interoperability between actors and placed full payloads on the ledger. Griggs et al.\cite{griggs2018healthcare} applied smart contracts to medical IoT data with HIPAA-compliant notifications, yet reported transaction delays that limit real-time utility. Quaini et al.\cite{quaini2018model} presented UniRec, a distributed electronic health record architecture whose full-replication design produces prohibitive disk space requirements at scale. These first-generation approaches share a common limitation: treating the blockchain as a bulk data repository rather than a coordination layer, thereby inheriting the storage, cost, and privacy weaknesses identified in Section~\ref{sec:introduction}.

More recent architectures have moved toward hybrid storage and privacy-layered designs. Zhang et al.\cite{zhang2018fhirchain} proposed FHIRChain, combining blockchain coordination with HL7 FHIR interoperability for collaborative cancer care, but reported semantic interoperability challenges and legacy system limitations. The architectural direction established by such hybrid approaches---separating data custody from on-chain coordination---informs the selective-anchoring model proposed in this work, which extends it with encrypted off-chain storage via IPFS, zero-knowledge verification, and dynamic permission lifecycle management.

\subsubsection{Surveillance-Specific Digital Architectures}

Digital surveillance systems have predominantly employed machine learning. Bollig et al.\cite{bollig2020machine} applied LSTM networks to syndromic surveillance across 33,000 veterinary necropsy reports, demonstrating temporal and spatial pattern detection. Harvey et al.\cite{harvey2021predicting} developed Gaussian Process and Random Forest models for malaria early warning in Burkina Faso using routine health facility data. Rubio-Solis et al.\cite{rubiosolis2019zika} proposed neural network-based predictive surveillance for dengue at district level in Brazil. While these systems demonstrate predictive utility, they require centralized data aggregation, remain vulnerable to adversarial manipulation, and provide no inherent guarantees of provenance, immutability, or audit traceability---properties essential for trustworthy multi-institutional surveillance.

\subsubsection{Architectural Gap Analysis}

Table~\ref{tab:related} positions the present work against prior systems across the six architectural dimensions identified in the blueprint.

\begin{table*}[!ht]
\caption{Comparative analysis of related work against target architectural dimensions.}\label{tab:related}
\begin{tabular*}{\textwidth}{@{\extracolsep\fill}llllll@{}}
\toprule
Study & Hybrid Storage & Layered Privacy & Edge/Fog & Consensus Rationale & LMIC Awareness \\
\midrule
Omar et al.\cite{omar2019privacy} & No & No & No & No & No \\
Griggs et al.\cite{griggs2018healthcare} & No & Partial & No & No & No \\
Quaini et al.\cite{quaini2018model} & No & No & No & No & No \\
Zhang et al.\cite{zhang2018fhirchain} & Partial & Partial & No & No & No \\
Bollig et al.\cite{bollig2020machine} & N/A & No & No & N/A & No \\
Harvey et al.\cite{harvey2021predicting} & N/A & No & No & N/A & Partial \\
Rubio-Solis et al.\cite{rubiosolis2019zika} & N/A & No & No & N/A & No \\
\textbf{This work} & \textbf{Yes} & \textbf{Yes} & \textbf{Yes} & \textbf{Yes} & \textbf{Yes} \\
\bottomrule
\end{tabular*}
\end{table*}

\textbf{Research gap.} No prior blockchain-health architecture integrates all five dimensions required for LMIC disease surveillance: (1)~hybrid storage with selective on-chain anchoring, (2)~layered privacy with zero-knowledge verification and dynamic permission lifecycle, (3)~edge/fog preprocessing for resource-constrained peripheral sites, (4)~comparative consensus evaluation with explicit validator governance, and (5)~LMIC-aware asymmetric node-role design aligned to public-health administrative hierarchies. The present work addresses this compound gap by proposing a unified architecture that treats privacy, consensus, scalability, and resource asymmetry as linked governance concerns within a single hierarchical surveillance framework.

%% ===== RESEARCH METHODOLOGY =====
\section{Research Methodology}\label{sec:methodology}

\subsection{Research Design}\label{subsec:design}

This study adopts the Systems Development Research Methodology (SDRM) proposed by Nunamaker et al.\cite{nunamaker1990systems}, instantiated through five iterative stages aligned to the hybrid surveillance architecture (Figure~\ref{fig:sdrm}).

\begin{figure}[!ht]
\centering
\fbox{\parbox{0.8\columnwidth}{\centering\vspace{2em}[SDRM Process Diagram]\\Figure placeholder\vspace{2em}}}
\caption{Systems Development Research Methodology (SDRM) process.}\label{fig:sdrm}
\end{figure}

\textbf{Stage 1---Problem Identification and Concept Design.} A structured literature review identified five persistent deficiencies in LMIC surveillance systems (Section~\ref{sec:introduction}) and four categories of architectural shortcoming in first-generation blockchain-health designs: monolithic on-chain storage, narrow privacy models, unargued consensus assumptions, and resource-uniform participation. These findings established the requirements for a hybrid, privacy-layered architecture.

\textbf{Stage 2---Architecture Development.} Six design objectives were derived from the gap analysis: (1)~selective on-chain anchoring with encrypted off-chain storage, (2)~layered privacy with zero-knowledge verification and dynamic permission lifecycle, (3)~edge/fog preprocessing for resource-constrained sites, (4)~comparative consensus evaluation with validator governance, (5)~hierarchical event processing aligned to public-health administrative tiers, and (6)~FHIR-aware interoperability workflows. Each objective maps to a specific architectural component (Section~\ref{sec:model}).

\textbf{Stage 3---Analysis and Design.} Smart-contract coordination logic, off-chain storage interfaces (IPFS), access-control state machines with grant/revoke/expire semantics, and tier-specific node-role specifications were designed against the six objectives.

\textbf{Stage 4---Prototype Development.} A functional prototype was implemented on Ethereum to validate the smart-contract coordination layer, demonstrating selective anchoring, hierarchical access control, and automated completeness verification.

\textbf{Stage 5---Observation and Evaluation.} The prototype was subjected to three-dimensional evaluation: performance (block generation timing), usability (Software Usability Measurement Inventory with domain professionals), and security (expanded threat model addressing validator collusion, metadata leakage, overbroad tier exposure, and equity-aware governance).

\subsection{Blockchain Suitability Assessment}\label{subsec:suitability}

Prior to architecture design, blockchain suitability for the DEWS domain was assessed using the IEEE Spectrum decision framework\cite{peck2017blockchain}, selected from eight published frameworks (Figure~\ref{fig:frameworks}) for its parsimony and accessibility to domain stakeholders.

\begin{figure}[!ht]
\centering
\fbox{\parbox{0.8\columnwidth}{\centering\vspace{2em}[Blockchain Decision Frameworks]\\Figure placeholder\vspace{2em}}}
\caption{Blockchain decision frameworks evaluated.}\label{fig:frameworks}
\end{figure}

Ten purposively sampled epidemiologists and decision-making staff at the Amhara Public Health Institute completed semi-structured interviews (February--March 2021). Table~\ref{tab:survey} summarizes the responses.

\begin{table}[!ht]
\caption{Blockchain suitability survey results (n=10).}\label{tab:survey}
\begin{tabular*}{\columnwidth}{@{\extracolsep\fill}lc@{}}
\toprule
Decision Criterion & Affirmative (\%) \\
\midrule
Traditional methods inadequate & 90 \\
Multiple participants need data updates & 100 \\
Low mutual trust among actors & 90 \\
Need for shared data store & 80 \\
Need for transaction transparency & 70 \\
Need for immutable records & 80 \\
Need for decentralized control & 70 \\
\midrule
Average & 82.85 \\
\bottomrule
\end{tabular*}
\end{table}

An average of 82.85\% affirmative responses exceeded the framework's majority threshold, confirming blockchain suitability for the surveillance domain. This assessment validates the use of blockchain \emph{as a coordination and proof layer}; it does not justify full on-chain data storage, which the architecture explicitly avoids. The specific architectural pattern---selective anchoring with encrypted off-chain storage, layered privacy, and comparative consensus---was derived from the literature-driven gap analysis rather than from the suitability survey alone.

\subsection{Scope and Limitations}\label{subsec:scope}

The study scope encompasses the design, prototyping, and evaluation of the smart-contract coordination layer within the hybrid architecture. The following methodological boundaries apply: (1)~the prototype validates coordination-layer feasibility on Ethereum but does not implement the complete off-chain storage, edge/fog, or Layer-2 components, which are specified architecturally; (2)~performance evaluation uses the Rinkeby test network, not mainnet deployment; (3)~suitability survey participants are drawn from a single regional institute (n=10), limiting generalizability of stakeholder findings; and (4)~the expanded threat model is analytically derived rather than empirically penetration-tested. These constraints bound the current empirical claims while preserving the architectural contribution, which is designed to be transferable across LMIC surveillance contexts.

%% ===== PROPOSED MODEL =====
\section{Proposed Hybrid Blockchain Architecture for Disease Early Warning and Surveillance}\label{sec:model}

This section presents a hybrid, privacy-layered, hierarchical architecture for infectious disease early warning and surveillance (DEWS) in low- and middle-income country (LMIC) settings. The architecture replaces monolithic on-chain storage with \emph{selective anchoring}: encrypted surveillance payloads reside off-chain while the Ethereum blockchain retains only cryptographic hashes, access-control state, provenance records, alerts, and audit events. This separation addresses four categories of shortcoming in first-generation blockchain-health systems: unsustainable storage replication, narrow immutability-only privacy, unargued consensus assumptions, and resource-uniform participation models that exclude constrained facilities.

\subsection{Multi-Tier Architectural Topology}\label{subsec:topology}

The architecture organizes participants into five functionally distinct tiers aligned with the Ethiopian public health administrative hierarchy. Figure~\ref{fig:architecture} illustrates the overall topology.

\begin{figure}[!ht]
\centering
\fbox{\parbox{0.8\columnwidth}{\centering\vspace{2em}[Multi-Tier Architectural Topology]\\Figure placeholder\vspace{2em}}}
\caption{Hybrid blockchain architecture for disease early warning and surveillance showing five processing tiers.}\label{fig:architecture}
\end{figure}

% ===== Mermaid Source =====
% Figure: Overall Architectural Model
% Scientific purpose: Provides a single high-level view of the five-tier hybrid architecture,
% showing how surveillance data flows from community capture through edge preprocessing
% and district batching to selective blockchain anchoring with off-chain encrypted storage.
%
% graph TB
%     subgraph T1["Tier 1 · Community / Peripheral Capture"]
%         CHW["Community Health Workers"]
%         HF["Health Facilities"]
%     end
%
%     subgraph T2["Tier 2 · Edge / Fog Processing"]
%         EFN["Edge/Fog Nodes"]
%         VAL["Preliminary Validation"]
%         BUF["Local Buffering & Batching"]
%     end
%
%     subgraph T3["Tier 3 · District / Regional Coordination"]
%         DC["District Coordinator"]
%         BA["Batch Aggregation"]
%         SV["Supervisory Validation"]
%     end
%
%     subgraph T4["Tier 4 · Blockchain Coordination"]
%         SC["Smart Contracts"]
%         HA["Hash Anchoring"]
%         ACL["Access Control & Permissions"]
%         CON["Consent & Provenance Records"]
%         AU["Audit Trails"]
%         AL["Alert Events"]
%     end
%
%     subgraph T5["Tier 5 · Off-Chain Data Layer"]
%         IPFS["IPFS / Encrypted Storage"]
%         ENC["AES-256 Encrypted Payloads"]
%     end
%
%     CHW -->|"encrypted event"| EFN
%     HF -->|"encrypted event"| EFN
%     EFN --> VAL
%     VAL --> BUF
%     BUF -->|"validated batch"| DC
%     DC --> BA
%     DC --> SV
%     BA -->|"batch tx"| SC
%     SC --> HA
%     SC --> ACL
%     SC --> CON
%     SC --> AU
%     SC --> AL
%     EFN -.->|"encrypted payload"| IPFS
%     DC -.->|"encrypted payload"| IPFS
%     IPFS --- ENC
%     HA -.->|"hash/pointer"| IPFS
%
%     style T1 fill:#f5f5f5,stroke:#333
%     style T2 fill:#e8e8e8,stroke:#333
%     style T3 fill:#dcdcdc,stroke:#333
%     style T4 fill:#d0d0d0,stroke:#333
%     style T5 fill:#c4c4c4,stroke:#333
% ==========================

\textbf{Tier~1: Community and Peripheral Capture Layer.} Community health workers and primary health posts originate surveillance events. These nodes operate as lightweight clients---they encrypt case data locally and submit payloads to off-chain storage without running validator infrastructure. This tier accommodates intermittent connectivity, underpowered devices, and limited staff capacity characteristic of LMIC peripheral sites.

\textbf{Tier~2: Edge/Fog Processing Layer.} Near-source processing nodes perform preliminary validation, data formatting, local buffering during connectivity interruptions, and transaction batching. This layer reduces bandwidth demands and latency for constrained facilities while ensuring that submissions conform to schema requirements before propagation.

\textbf{Tier~3: District/Regional Coordination Layer.} Woreda, zone, and regional health offices serve as intermediate trust and aggregation nodes. These facilities batch validated submissions from subordinate tiers, perform supervisory review, and prepare aggregate transaction sets for blockchain anchoring. The coordination layer mirrors the existing public health reporting hierarchy, enabling bidirectional information flow---upward for case reporting and downward for supervisory feedback.

\textbf{Tier~4: Blockchain Coordination Layer.} Ethereum smart contracts provide the system's coordination, enforcement, provenance, and auditability functions. On-chain data is strictly limited to: cryptographic hashes of off-chain payloads, content-addressable pointers, permission and access-control state, consent and provenance records, alert events, and audit trail entries. Raw surveillance payloads are never stored on-chain.

\textbf{Tier~5: Off-Chain Encrypted Data Layer.} Bulk surveillance records---case reports, laboratory results, specimen metadata---reside in encrypted off-chain storage using IPFS content addressing. Each stored object is linked to its on-chain hash anchor, enabling integrity verification without ledger replication of sensitive health data.

\subsection{Selective Anchoring and Core Data Flow}\label{subsec:dataflow}

The selective anchoring principle governs all data movement through the architecture. Figure~\ref{fig:smartcontract} illustrates the internal smart contract logic, while Figure~\ref{fig:web3} presents the end-to-end data flow sequence. The workflow proceeds as follows:

\begin{figure}[!ht]
\centering
\fbox{\parbox{0.8\columnwidth}{\centering\vspace{2em}[Selective Anchoring Data Flow]\\Figure placeholder\vspace{2em}}}
\caption{Selective anchoring workflow showing separation of encrypted off-chain storage from on-chain coordination artifacts.}\label{fig:smartcontract}
\end{figure}

% ===== Mermaid Source =====
% Figure: Smart Contract Selective-Anchoring Mechanism
% Scientific purpose: Illustrates the internal logic of the smart contract layer,
% showing which artifacts are anchored on-chain (hashes, permissions, alerts, audit)
% versus which remain off-chain (encrypted payloads), enforcing the selective-anchoring principle.
%
% flowchart LR
%     subgraph Input["Incoming Batch from District"]
%         BT["Batch Transaction"]
%         HP["Hash + Pointer"]
%         PM["Permission Metadata"]
%     end
%
%     subgraph SC["Smart Contract Logic"]
%         VER["Verify Batch Integrity"]
%         ANC["Anchor Hash On-Chain"]
%         PERM["Update ACL State"]
%         PROV["Record Consent & Provenance"]
%         ALERT["Evaluate Alert Rules"]
%         AUDIT["Write Audit Event"]
%     end
%
%     subgraph OnChain["On-Chain State"]
%         H["Hashes & Pointers"]
%         P["Permissions & Consent"]
%         PR["Provenance Records"]
%         A["Alert Records"]
%         T["Audit Trail"]
%     end
%
%     subgraph OffChain["Off-Chain"]
%         IPFS["Encrypted Payloads · IPFS"]
%     end
%
%     BT --> VER
%     HP --> VER
%     PM --> VER
%     VER --> ANC
%     VER --> PERM
%     VER --> PROV
%     VER --> ALERT
%     VER --> AUDIT
%     ANC --> H
%     PERM --> P
%     PROV --> PR
%     ALERT --> A
%     AUDIT --> T
%     H -.->|"pointer reference"| IPFS
%
%     style OnChain fill:#e8e8e8,stroke:#333
%     style OffChain fill:#f5f5f5,stroke:#333,stroke-dasharray:5 5
% ==========================


\begin{figure}[!ht]
\centering
\fbox{\parbox{0.8\columnwidth}{\centering\vspace{2em}[End-to-End Data Flow]\\Figure placeholder\vspace{2em}}}
\caption{End-to-end data flow from event origination through encryption, off-chain storage, hash anchoring, and authorized retrieval. Solid arrows represent payload movement; dashed arrows represent hash/pointer references.}\label{fig:web3}
\end{figure}

% ===== Mermaid Source =====
% Figure: Core Data Flow
% Scientific purpose: Replaces the original Web3.js-centric figure with a complete
% end-to-end data flow diagram showing encryption-first design, off-chain storage,
% selective anchoring, and authorized retrieval—the central data lifecycle of the
% revised architecture.
%
% sequenceDiagram
%     participant RA as Reporting Actor
%     participant EF as Edge/Fog Node
%     participant DC as District Coordinator
%     participant OC as Off-Chain Store (IPFS)
%     participant BC as Blockchain (Smart Contract)
%     participant AU as Authorized Verifier
%
%     RA->>EF: Submit encrypted surveillance event
%     EF->>EF: Validate, format, buffer
%     EF->>OC: Store encrypted payload
%     OC-->>EF: Return content-address (CID)
%     EF->>DC: Forward validated event + CID
%     DC->>DC: Batch and aggregate
%     DC->>BC: Anchor batch (hashes, CIDs, permissions)
%     BC-->>DC: Confirm anchoring tx
%     Note over BC: On-chain: hash, pointer, ACL, audit
%     AU->>BC: Request access + verify provenance
%     BC-->>AU: Return CID + permission token
%     AU->>OC: Retrieve encrypted payload by CID
%     OC-->>AU: Return encrypted payload
%     AU->>AU: Decrypt with authorized key
% ==========================

\begin{enumerate}
\item A reporting actor at Tier~1 creates a surveillance event (case report, laboratory result, or completeness submission).
\item The payload is encrypted using the actor's cryptographic credentials before leaving the originating node.
\item The encrypted payload is stored off-chain via IPFS, yielding a content-addressed hash.
\item The hash, a storage pointer, permission metadata, and relevant audit information are prepared for blockchain anchoring.
\item Edge/fog nodes (Tier~2) or district/regional coordinators (Tier~3) preprocess, validate, and batch submissions.
\item Smart contracts at Tier~4 commit only the selective on-chain artifacts---hash, permissions, provenance, and alert state.
\item Authorized parties verify data integrity, provenance, and access status through on-chain proofs and permission logic, retrieving encrypted payloads from off-chain storage only when access is granted.
\item Alert events and audit records are preserved on-chain as immutable traceability artifacts.
\end{enumerate}

This workflow eliminates full ledger replication of surveillance data, reduces gas costs and storage growth, and ensures that bulk health records are never directly exposed on the distributed ledger.

\subsection{Hierarchical Disease Reporting}\label{subsec:reporting}

The disease reporting component implements Ethiopia's five-tier public health surveillance hierarchy: community health posts, woreda health offices, zone health offices, regional health offices, and the EHNRI/PHEM national center.

The revised architecture separates the reporting workflow into four functional phases:

\textbf{Capture phase.} Community and facility reporters (Tier~1) encrypt case data and submit encrypted payloads to IPFS. Each submission generates a content hash and metadata record. Reporters operate as lightweight clients with no blockchain validation responsibilities.

\textbf{Aggregation phase.} Edge/fog nodes (Tier~2) buffer, validate, and batch submissions from multiple reporters. District and regional coordinators (Tier~3) aggregate validated batches and perform hierarchical completeness checks before anchoring.

\textbf{Coordination phase.} Smart contracts at Tier~4 receive batch anchoring transactions containing hashes, permission grants, and provenance records. A factory contract instantiates and manages organizational report-tracking contracts. Access control enforces hierarchical constraints: higher-level offices can retrieve and review reports from subordinate organizations; subordinate levels cannot access peer-level or superior-level reports unless explicitly authorized through on-chain permission grants.

\textbf{Audit phase.} Every report submission, review action, feedback message, and permission change constitutes a permanent on-chain audit record. Supervisory feedback flows downward through the same contract interface, maintaining bidirectional traceability without storing report content on the ledger.

Figure~\ref{fig:reporting} illustrates the redesigned reporting process.

\begin{figure}[!ht]
\centering
\fbox{\parbox{0.8\columnwidth}{\centering\vspace{2em}[Hierarchical Disease Reporting Process]\\Figure placeholder\vspace{2em}}}
\caption{Hierarchical disease reporting process with selective anchoring and role-asymmetric participation.}\label{fig:reporting}
\end{figure}

% ===== Mermaid Source =====
% Figure: Hierarchical Disease Reporting Process
% Scientific purpose: Shows the five-tier Ethiopian public health reporting hierarchy
% with the revised data handling: encrypted off-chain storage at each tier,
% selective hash anchoring, and bidirectional flow (upward reports, downward feedback).
%
% graph TB
%     CHP["Community Health Post"]
%     WHO["Woreda Health Office"]
%     ZHO["Zone Health Office"]
%     RHO["Regional Health Office"]
%     EHNRI["EHNRI/PHEM Center"]
%     OC["Off-Chain Encrypted Store"]
%     BC["Blockchain · Hash Anchor"]
%
%     CHP -->|"encrypted report ↑"| WHO
%     WHO -->|"aggregated batch ↑"| ZHO
%     ZHO -->|"aggregated batch ↑"| RHO
%     RHO -->|"aggregated batch ↑"| EHNRI
%
%     EHNRI -.->|"feedback ↓"| RHO
%     RHO -.->|"feedback ↓"| ZHO
%     ZHO -.->|"feedback ↓"| WHO
%     WHO -.->|"feedback ↓"| CHP
%
%     CHP -->|"payload"| OC
%     WHO -->|"payload"| OC
%     ZHO -->|"payload"| OC
%     RHO -->|"payload"| OC
%
%     WHO -->|"hash + audit"| BC
%     ZHO -->|"hash + audit"| BC
%     RHO -->|"hash + audit"| BC
%     EHNRI -->|"hash + audit"| BC
%
%     style OC fill:#e8e8e8,stroke:#333,stroke-dasharray:5 5
%     style BC fill:#d0d0d0,stroke:#333
% ==========================

\subsection{Laboratory Network Integration}\label{subsec:laboratory}

The laboratory component addresses pathogen identification and result exchange across the four-tiered laboratory network: woreda, zone, regional, and national reference levels.

Under the hybrid architecture, laboratory results and specimen metadata are encrypted and stored off-chain. On-chain contracts manage only provenance records, result-confirmation events, and access permissions. Forward (upward) reporting transmits hash-anchored identification results to higher-level facilities for confirmation. Backward (downward) communication enables reference laboratories to disseminate confirmations, corrections, and follow-up information through permission-controlled off-chain retrieval with on-chain audit logging.

Selective disclosure ensures that specimen-level data is accessible only to authorized laboratory tiers. Dynamic access grants allow temporary cross-tier visibility during outbreak investigations, with automatic expiry preventing persistent overbroad access. Figure~\ref{fig:laboratory} depicts the laboratory network design.

\begin{figure}[!ht]
\centering
\fbox{\parbox{0.8\columnwidth}{\centering\vspace{2em}[Laboratory Network Design]\\Figure placeholder\vspace{2em}}}
\caption{Laboratory network integration with off-chain encrypted storage and on-chain provenance anchoring.}\label{fig:laboratory}
\end{figure}

% ===== Mermaid Source =====
% Figure: Laboratory Network Reporting Design
% Scientific purpose: Depicts the four-tier laboratory network with encrypted off-chain
% specimen storage and blockchain hash anchoring, showing bidirectional flow for
% specimen submission (upward) and result confirmation (downward).
%
% graph TB
%     WL["Woreda Laboratory"]
%     ZL["Zone Laboratory"]
%     RL["Regional Laboratory"]
%     NL["National Reference Lab"]
%     OC["Off-Chain Encrypted Store"]
%     BC["Blockchain · Provenance Anchor"]
%
%     WL -->|"encrypted specimen data ↑"| ZL
%     ZL -->|"encrypted specimen data ↑"| RL
%     RL -->|"encrypted specimen data ↑"| NL
%
%     NL -.->|"confirmation ↓"| RL
%     RL -.->|"confirmation ↓"| ZL
%     ZL -.->|"confirmation ↓"| WL
%
%     WL -->|"payload"| OC
%     ZL -->|"payload"| OC
%     RL -->|"payload"| OC
%
%     ZL -->|"hash + provenance"| BC
%     RL -->|"hash + provenance"| BC
%     NL -->|"hash + provenance"| BC
%
%     style OC fill:#e8e8e8,stroke:#333,stroke-dasharray:5 5
%     style BC fill:#d0d0d0,stroke:#333
% ==========================

\subsection{Layered Privacy Architecture}\label{subsec:privacy}

The privacy model replaces immutability-only protection with three complementary mechanisms:

\textbf{Encrypted off-chain storage.} All surveillance payloads are encrypted before storage, ensuring that compromise of the storage layer alone does not expose case-level data.

\textbf{Dynamic access control.} Smart contracts implement fine-grained permission management with grant, revoke, and time-bound expire semantics. Access rights are scoped to specific organizational tiers and reporting relationships, preventing overbroad cross-tier exposure. Permission state changes are recorded on-chain for auditability.

\textbf{Zero-knowledge verification.} Privacy-preserving validation enables authorized parties to verify properties of surveillance submissions---such as completeness, timeliness, or demographic eligibility---without accessing underlying sensitive records. Zero-knowledge proofs support selective disclosure during cross-organizational coordination and outbreak response while minimizing data exposure.

Figure~\ref{fig:privacy} illustrates the layered privacy control architecture.

\begin{figure}[!ht]
\centering
\fbox{\parbox{0.8\columnwidth}{\centering\vspace{2em}[Layered Privacy Architecture]\\Figure placeholder\vspace{2em}}}
\caption{Layered privacy architecture with zero-knowledge verification. The three privacy mechanisms---encrypted off-chain storage, dynamic access control with revocation, and zero-knowledge proof verification---operate across architectural tiers to enforce selective disclosure without exposing raw surveillance records.}\label{fig:privacy}
\end{figure}

% ===== Mermaid Source =====
% Figure: Privacy and Zero-Knowledge Verification Layer
% Scientific purpose: Visualizes the three-mechanism privacy architecture that replaces
% the thesis's immutability-only privacy model. Shows how encrypted storage, dynamic ACL,
% and ZKP work together across tiers to enable selective disclosure and revocable access.
%
% graph TB
%     subgraph PL["Privacy Layer"]
%         direction TB
%         subgraph M1["Mechanism 1 · Encrypted Off-Chain Storage"]
%             ENC["AES-256 Encryption"]
%             IPFS["IPFS Content-Addressed Store"]
%         end
%         subgraph M2["Mechanism 2 · Dynamic Access Control"]
%             GRANT["Grant Permission"]
%             REVOKE["Revoke Permission"]
%             EXPIRE["Time-Bound Expiry"]
%             ACL["Smart Contract ACL"]
%         end
%         subgraph M3["Mechanism 3 · Zero-Knowledge Verification"]
%             PROVE["Prover (Reporting Tier)"]
%             VERIFY["Verifier (Requesting Tier)"]
%             ZKP["zk-SNARK Circuit"]
%         end
%     end
%
%     ENC --> IPFS
%     GRANT --> ACL
%     REVOKE --> ACL
%     EXPIRE --> ACL
%     PROVE -->|"proof π"| ZKP
%     ZKP -->|"accept/reject"| VERIFY
%
%     style M1 fill:#f5f5f5,stroke:#333
%     style M2 fill:#e8e8e8,stroke:#333
%     style M3 fill:#dcdcdc,stroke:#333
% ==========================

\subsection{Report Completeness Verification}\label{subsec:completeness}

The Ethiopian DEWS protocol mandates weekly reporting; a report is classified as complete when all reporting units within a designated catchment area have submitted within the prescribed period.

Under the hybrid architecture, completeness verification operates on anchored events rather than full ledger queries. Smart contracts at each hierarchical level maintain a registry of expected reporters and track submission-anchor events. The completeness calculator function computes the ratio of received anchoring events to expected submissions for a given period and catchment area. Because only hashes and metadata are queried---not full surveillance payloads---the verification mechanism scales efficiently without requiring ledger-wide data replication.

Edge/fog nodes assist completeness tracking by maintaining local submission registries and alerting coordinators to missing reports before the reporting deadline, enabling proactive gap identification. Figure~\ref{fig:completeness} illustrates the event-driven completeness verification design.

\begin{figure}[!ht]
\centering
\fbox{\parbox{0.8\columnwidth}{\centering\vspace{2em}[Completeness Verification]\\Figure placeholder\vspace{2em}}}
\caption{Event-driven report completeness verification using anchored submission events.}\label{fig:completeness}
\end{figure}

% ===== Mermaid Source =====
% Figure: Report Completeness Verification
% Scientific purpose: Shows how completeness is verified using on-chain hash records
% rather than raw data, aligning with the selective-anchoring principle while
% automating the weekly reporting compliance check.
%
% flowchart TD
%     REG["Reporting-Period Registry"]
%     SC["Completeness Calculator Contract"]
%     HA["On-Chain Hash Index"]
%     EF["Edge/Fog Local Registry"]
%     CR["Completeness Ratio Output"]
%     ALERT["Gap Alert Event"]
%
%     REG -->|"expected units & period"| SC
%     HA -->|"submitted hashes per unit"| SC
%     EF -->|"local submission status"| SC
%     SC -->|"received / expected"| CR
%     SC -->|"missing units"| ALERT
%     EF -.->|"proactive gap alert"| ALERT
%
%     style SC fill:#e8e8e8,stroke:#333
%     style HA fill:#d0d0d0,stroke:#333
%     style EF fill:#f0f0f0,stroke:#333,stroke-dasharray:3 3
% ==========================

\subsection{Surveillance Data Analysis}\label{subsec:analysis}

The data analysis component implements three epidemiological dimensions: temporal (day, week, month), spatial (kebele, woreda, town), and person-based (age, sex, occupation) analysis for outbreak detection and disease burden assessment. Figures~\ref{fig:temporal},~\ref{fig:spatial}, and~\ref{fig:person} detail the respective analysis workflows.

Under the revised architecture, analytics are separated from ledger settlement. Computationally intensive epidemiological analysis---including incidence proportion and prevalence calculations---executes at the district/regional coordination layer (Tier~3) over decrypted off-chain data accessible through permission-controlled retrieval. Only summary indicators, alert triggers, and anomaly flags are anchored on-chain as audit-worthy events.

This placement ensures that: (1) analysis operates on complete decrypted datasets rather than on-chain transaction data; (2) privacy-preserving aggregation can be applied before results cross organizational boundaries; (3) constrained peripheral nodes are not burdened with analytical computation; and (4) early warning alerts inherit the integrity guarantees of blockchain anchoring without requiring raw case data on the ledger.


\begin{figure}[!ht]
\centering
\fbox{\parbox{0.8\columnwidth}{\centering\vspace{2em}[Temporal Analysis Workflow]\\Figure placeholder\vspace{2em}}}
\caption{Temporal analysis workflow. Authorized analysts retrieve encrypted case data from off-chain storage, verify integrity against on-chain hashes, decrypt, and compute time-series trend indicators for outbreak detection.}\label{fig:temporal}
\end{figure}

% ===== Mermaid Source =====
% Figure: Temporal Analysis Workflow
% Scientific purpose: Shows how temporal epidemiological analysis operates on
% integrity-verified off-chain data, maintaining the privacy-layered architecture
% while computing time-series outbreak indicators.
%
% flowchart LR
%     BC["Blockchain · Hash Anchor"]
%     OC["Off-Chain Encrypted Store"]
%     AN["Authorized Analyst"]
%     DEC["Decrypt & Parse"]
%     TS["Time-Series Extraction"]
%     TR["Trend & Anomaly Detection"]
%     OUT["Outbreak Alert Output"]
%
%     AN -->|"verify provenance"| BC
%     BC -.->|"CID + permission"| AN
%     AN -->|"retrieve by CID"| OC
%     OC -->|"encrypted payload"| DEC
%     DEC --> TS
%     TS --> TR
%     TR --> OUT
%
%     style BC fill:#d0d0d0,stroke:#333
%     style OC fill:#e8e8e8,stroke:#333,stroke-dasharray:5 5
% ==========================


\begin{figure}[!ht]
\centering
\fbox{\parbox{0.8\columnwidth}{\centering\vspace{2em}[Spatial Analysis Workflow]\\Figure placeholder\vspace{2em}}}
\caption{Spatial analysis workflow. Geospatial variables are extracted from integrity-verified off-chain records to identify geographic clustering and track disease spread across administrative boundaries.}\label{fig:spatial}
\end{figure}

% ===== Mermaid Source =====
% Figure: Spatial Analysis Workflow
% Scientific purpose: Illustrates extraction of geospatial variables from verified
% off-chain records for areal spread analysis across hierarchical administrative units.
%
% flowchart LR
%     BC["Blockchain · Hash Anchor"]
%     OC["Off-Chain Encrypted Store"]
%     AN["Authorized Analyst"]
%     DEC["Decrypt & Parse"]
%     GEO["Geospatial Extraction"]
%     CL["Cluster & Spread Analysis"]
%     MAP["Spatial Alert Output"]
%
%     AN -->|"verify provenance"| BC
%     BC -.->|"CID + permission"| AN
%     AN -->|"retrieve by CID"| OC
%     OC -->|"encrypted payload"| DEC
%     DEC --> GEO
%     GEO --> CL
%     CL --> MAP
%
%     style BC fill:#d0d0d0,stroke:#333
%     style OC fill:#e8e8e8,stroke:#333,stroke-dasharray:5 5
% ==========================


\begin{figure}[!ht]
\centering
\fbox{\parbox{0.8\columnwidth}{\centering\vspace{2em}[Person Analysis Workflow]\\Figure placeholder\vspace{2em}}}
\caption{Person-based analysis workflow. Demographic variables are extracted from integrity-verified off-chain records to identify disproportionately affected population subgroups while preserving individual-level privacy through selective disclosure.}\label{fig:person}
\end{figure}

% ===== Mermaid Source =====
% Figure: Person-Based Analysis Workflow
% Scientific purpose: Shows demographic analysis operating on verified off-chain data
% with privacy-preserving selective disclosure, enabling subgroup identification
% without exposing individual-level records unnecessarily.
%
% flowchart LR
%     BC["Blockchain · Hash Anchor"]
%     OC["Off-Chain Encrypted Store"]
%     AN["Authorized Analyst"]
%     SD["Selective Disclosure Filter"]
%     DEM["Demographic Extraction"]
%     CMP["Subgroup Comparison"]
%     OUT["Population Risk Output"]
%
%     AN -->|"verify provenance"| BC
%     BC -.->|"CID + permission"| AN
%     AN -->|"retrieve by CID"| OC
%     OC -->|"encrypted payload"| SD
%     SD -->|"filtered fields"| DEM
%     DEM --> CMP
%     CMP --> OUT
%
%     style BC fill:#d0d0d0,stroke:#333
%     style OC fill:#e8e8e8,stroke:#333,stroke-dasharray:5 5
% ==========================


\begin{figure}[!ht]
\centering
\fbox{\parbox{0.8\columnwidth}{\centering\vspace{2em}[Incidence Analysis Workflow]\\Figure placeholder\vspace{2em}}}
\caption{Incidence proportion analysis workflow. New-case counts are derived from integrity-verified off-chain records and divided by the at-risk population denominator to produce incidence rates for outbreak detection.}\label{fig:incidence}
\end{figure}

% ===== Mermaid Source =====
% Figure: Incidence Proportion Analysis Workflow
% Scientific purpose: Details the computation of incidence proportion from verified
% off-chain data, showing the numerator (new cases) and denominator (at-risk population)
% derivation pathway.
%
% flowchart TD
%     VD["Verified Decrypted Records"]
%     NC["New Case Count (numerator)"]
%     POP["At-Risk Population (denominator)"]
%     INC["Incidence Proportion = NC / POP"]
%     EW["Early Warning Threshold Check"]
%
%     VD --> NC
%     VD --> POP
%     NC --> INC
%     POP --> INC
%     INC --> EW
%
%     style INC fill:#e8e8e8,stroke:#333
% ==========================


\begin{figure}[!ht]
\centering
\fbox{\parbox{0.8\columnwidth}{\centering\vspace{2em}[Prevalence Analysis Workflow]\\Figure placeholder\vspace{2em}}}
\caption{Prevalence analysis workflow. Total existing and new case counts are derived from integrity-verified off-chain records and divided by total population to produce point or period prevalence measures for disease burden assessment.}\label{fig:prevalence}
\end{figure}

% ===== Mermaid Source =====
% Figure: Prevalence Analysis Workflow
% Scientific purpose: Details the computation of prevalence from verified off-chain data,
% showing the summation of new and pre-existing cases relative to total population.
%
% flowchart TD
%     VD["Verified Decrypted Records"]
%     EC["Existing + New Cases (numerator)"]
%     TPOP["Total Population (denominator)"]
%     PREV["Prevalence = EC / TPOP"]
%     DB["Disease Burden Assessment"]
%
%     VD --> EC
%     VD --> TPOP
%     EC --> PREV
%     TPOP --> PREV
%     PREV --> DB
%
%     style PREV fill:#e8e8e8,stroke:#333
% ==========================

\textbf{Morbidity measures.} Incidence proportion (new cases relative to the at-risk population, Figure~\ref{fig:incidence}) and prevalence (existing cases at a point or period, Figure~\ref{fig:prevalence}) are computed at Tier~3 coordination nodes.\cite{cdc2012principles} Results are anchored on-chain as verifiable indicator events, supporting outbreak detection decisions grounded in tamper-evident, provenance-tracked epidemiological evidence.

\subsection{Consensus Design and Validator Governance}\label{subsec:consensus}

The architecture requires explicit consensus-selection rationale rather than default adoption of Proof of Authority (PoA). Figure~\ref{fig:consensus} illustrates the evaluation framework and tiered validator governance model. The consensus design addresses three dimensions:

\textbf{Comparative evaluation.} PoA, Delegated Proof of Stake (DPoS), PBFT variants, and Raft-based alliance models are evaluated against: (1) the hierarchical trust topology of the Ethiopian health system, (2) LMIC infrastructure constraints including latency and node capacity, (3) transaction throughput requirements for batched surveillance submissions, and (4) fault tolerance under realistic connectivity assumptions.

\textbf{Validator appointment.} Validators are drawn from institutional actors at Tier~3 and above---district, regional, and national health authorities with stable infrastructure and institutional accountability. Community and peripheral facilities are explicitly excluded from validator responsibilities, consistent with the resource-aware role separation principle.

\textbf{Governance and resilience.} The architecture specifies collusion thresholds, failure assumptions, and trust concentration bounds. Validator appointment follows institutional authority rather than computational capacity, and governance procedures for validator addition, removal, and dispute resolution are defined to prevent capture by any single administrative tier.


\begin{figure}[!ht]
\centering
\fbox{\parbox{0.8\columnwidth}{\centering\vspace{2em}[Consensus and Validator Governance]\\Figure placeholder\vspace{2em}}}
\caption{Consensus selection and validator governance model. Validators are appointed across administrative tiers with asymmetric responsibilities. The consensus mechanism is selected through comparative evaluation against workload, trust, latency, and infrastructure constraints specific to LMIC deployment.}\label{fig:consensus}
\end{figure}

% ===== Mermaid Source =====
% Figure: Consensus and Validator Governance Model
% Scientific purpose: Addresses the blueprint requirement for explicit consensus justification
% and validator governance. Shows the comparative evaluation framework and the tiered
% validator appointment model with collusion/failure assumptions.
%
% graph TB
%     subgraph VG["Validator Governance"]
%         direction TB
%         NA["National Validators · Full nodes"]
%         RA["Regional Validators · Full nodes"]
%         DA["District Validators · Light nodes"]
%         FA["Facility Actors · Reporting only"]
%     end
%
%     subgraph CE["Consensus Evaluation Framework"]
%         direction LR
%         POA["PoA"]
%         DPOS["DPoS"]
%         PBFT["PBFT Variants"]
%         RAFT["Raft-Based Alliance"]
%     end
%
%     subgraph CR["Selection Criteria"]
%         WL["Workload Profile"]
%         TC["Trust Concentration"]
%         LAT["Latency Tolerance"]
%         INF["Infrastructure Capacity"]
%     end
%
%     WL --> CE
%     TC --> CE
%     LAT --> CE
%     INF --> CE
%
%     CE -->|"selected mechanism"| VG
%
%     NA -->|"appoints"| RA
%     RA -->|"appoints"| DA
%
%     style VG fill:#e8e8e8,stroke:#333
%     style CE fill:#f5f5f5,stroke:#333
%     style CR fill:#dcdcdc,stroke:#333
% ==========================

\subsection{Scalability Extensions}\label{subsec:scalability}

The architecture supports three composable scalability mechanisms as design modules:

\textbf{Layer-2 batch settlement.} District/regional coordinators aggregate validated submissions off-chain and submit batch proofs or commitments to Layer-1, reducing main-chain transaction frequency and gas costs proportionally to batch size.

\textbf{Hierarchical state partitioning.} Surveillance state is partitioned by administrative geography, allowing regional subsets of the anchoring ledger to operate semi-independently while maintaining cross-partition integrity verification at the national level.

\textbf{FHIR-aware interoperability.} Privacy workflows are designed for compatibility with Fast Healthcare Interoperability Resources (FHIR) standards, enabling selective disclosure across organizational boundaries while maintaining standards-compliant data exchange and audit semantics.



\subsection{Threat-Response Mapping}\label{subsec:threats}

The architecture systematically addresses the four threat categories identified in the critique: confidentiality threats (metadata exposure, case-level overdisclosure), privacy threats (unauthorized access persistence, validation without confidentiality), governance threats (validator collusion/failure, trust concentration), and inclusion threats (heavy-node exclusion, integration weakness). Figure~\ref{fig:threats} maps each threat to its architectural countermeasure.

\begin{figure}[!ht]
\centering
\fbox{\parbox{0.8\columnwidth}{\centering\vspace{2em}[Threat-Response Mapping]\\Figure placeholder\vspace{2em}}}
\caption{Architectural threat-response mapping. Each identified threat category is linked to its corresponding architectural countermeasure, showing how the hybrid architecture systematically addresses confidentiality, privacy, governance, and inclusion threats identified in the critique.}\label{fig:threats}
\end{figure}

% ===== Mermaid Source =====
% Figure: Threat-Response Mapping
% Scientific purpose: Provides a systematic visualization of the expanded threat model,
% mapping each critique-identified threat to its architectural countermeasure. Ensures
% traceability between threats and design decisions across four threat categories.
%
% graph LR
%     subgraph CT["Confidentiality Threats"]
%         T1["Metadata Exposure"]
%         T2["Case-Level Overdisclosure"]
%     end
%     subgraph PT["Privacy Threats"]
%         T3["Unauthorized Access Persistence"]
%         T4["Validation Without Confidentiality"]
%     end
%     subgraph GT["Governance Threats"]
%         T5["Validator Collusion/Failure"]
%         T6["Trust Concentration"]
%     end
%     subgraph IT["Inclusion Threats"]
%         T7["Heavy-Node Exclusion"]
%         T8["Integration Weakness"]
%     end
%
%     subgraph RM["Architectural Responses"]
%         R1["Off-Chain Encrypted Storage"]
%         R2["Selective Disclosure + ZKP"]
%         R3["Dynamic ACL · Grant/Revoke/Expire"]
%         R4["Validator Governance Model"]
%         R5["Edge/Fog + Light Node Roles"]
%         R6["FHIR-Aware Privacy Workflows"]
%     end
%
%     T1 --> R1
%     T2 --> R1
%     T2 --> R2
%     T3 --> R3
%     T4 --> R2
%     T5 --> R4
%     T6 --> R4
%     T7 --> R5
%     T8 --> R6
%
%     style CT fill:#f5f5f5,stroke:#333
%     style PT fill:#ebebeb,stroke:#333
%     style GT fill:#e0e0e0,stroke:#333
%     style IT fill:#d6d6d6,stroke:#333
%     style RM fill:#cccccc,stroke:#333
% ==========================


%% ===== PROTOTYPE IMPLEMENTATION =====
\section{Prototype Implementation}\label{sec:implementation}

This section describes how a functional prototype realizes the five-tier hybrid architecture presented in Section~\ref{sec:model}. The prototype validates the smart-contract coordination layer (Tier~4) and demonstrates end-to-end data flow across community capture (Tier~1), off-chain encrypted storage (Tier~5), selective anchoring, hierarchical access control, and provenance tracking. Edge/fog preprocessing (Tier~2) and district-level batching (Tier~3) are implemented as application-layer modules; Layer-2 rollup and zero-knowledge verification components are specified architecturally and remain targets for future engineering. Figure~\ref{fig:proto_arch} presents the prototype software architecture mapped to the five-tier topology.

\begin{figure}[!ht]
\centering
\fbox{\parbox{0.8\columnwidth}{\centering\vspace{2em}[Prototype Software Architecture]\\Figure placeholder\vspace{2em}}}
\caption{Prototype software architecture mapped to the five-tier hybrid topology. Solid modules are implemented; dashed modules are architecturally specified.}\label{fig:proto_arch}
\end{figure}

% ===== Mermaid Source =====
% Figure: Prototype Software Architecture
% Scientific purpose: Maps the implemented prototype modules to the five architectural
% tiers, distinguishing fully implemented components from architecturally specified ones.
%
% graph TB
%     subgraph T1["Tier 1 · Community Capture Module"]
%         UI["React/Next.js Client"]
%         ENC_C["Client-Side AES-256 Encryption"]
%         MM["MetaMask Wallet"]
%     end
%
%     subgraph T2["Tier 2 · Edge/Fog Module"]
%         VAL_M["Validation Middleware"]
%         BUF_M["Local Buffer Queue"]
%     end
%
%     subgraph T3["Tier 3 · District Coordination Module"]
%         BATCH["Batch Aggregator"]
%         SUP["Supervisory Review Interface"]
%     end
%
%     subgraph T4["Tier 4 · Smart Contract Layer"]
%         RF["ReportFactory"]
%         RC["ReportCoordinator"]
%         ACL_C["AccessControl"]
%         COMP["CompletenessVerifier"]
%     end
%
%     subgraph T5["Tier 5 · Off-Chain Storage"]
%         IPFS_N["IPFS Node (ipfs-http-client)"]
%         ENC_S["Encrypted Payload Store"]
%     end
%
%     UI -->|"encrypted payload"| ENC_C
%     ENC_C -->|"ciphertext"| VAL_M
%     VAL_M --> BUF_M
%     BUF_M -->|"validated events"| BATCH
%     BATCH -->|"anchor tx"| RF
%     RF --> RC
%     RF --> ACL_C
%     RF --> COMP
%     ENC_C -.->|"encrypted blob"| IPFS_N
%     IPFS_N --- ENC_S
%     RC -.->|"CID pointer"| IPFS_N
%     MM -->|"sign tx"| RF
%     SUP -->|"review/rework"| RC
%
%     style T1 fill:#f5f5f5,stroke:#333
%     style T2 fill:#e8e8e8,stroke:#333
%     style T3 fill:#dcdcdc,stroke:#333
%     style T4 fill:#d0d0d0,stroke:#333
%     style T5 fill:#c4c4c4,stroke:#333
% =========================

\subsection{Development Environment and Technology Stack}\label{subsec:devenv}

The prototype was implemented using the following technology stack: Solidity~0.8 for smart contract development targeting the Ethereum Virtual Machine; Node.js as the server-side runtime; React and Next.js for the front-end application layer with server-side rendering; Web3.js for blockchain interaction; \texttt{ipfs-http-client} for IPFS content-addressed storage; MetaMask for client-side transaction signing; Ganache~CLI for local blockchain emulation during development; Mocha for automated contract testing; and Semantic~UI for interface theming. Client-side encryption uses the Web Crypto API with AES-256-GCM. Smart contracts were compiled using \texttt{solc} and deployed via a custom deployment script using the Truffle~HD Wallet Provider for account authentication.

The project comprises six modules: (1)~Ethereum components (smart contracts, compilation, deployment, and migration scripts); (2)~IPFS integration (encryption, upload, retrieval, and content-hash management); (3)~React page components implementing Tier~1 capture interfaces; (4)~Next.js routing, server-side rendering, and API middleware implementing Tier~2 validation logic; (5)~supervisory and batch coordination interfaces implementing Tier~3 workflows; and (6)~test components for contract and integration verification.

\subsection{Tier-by-Tier Implementation}\label{subsec:tiers}

\textbf{Tier~1: Community Data Capture.} Community health workers and facility reporters interact through a React-based web client. Upon case report submission, the client encrypts the surveillance payload using AES-256-GCM before any network transmission. The encrypted ciphertext is uploaded to IPFS via \texttt{ipfs-http-client}, which returns a content identifier (CID). The client then initiates a blockchain transaction through MetaMask, passing only the CID, the payload hash, and permission metadata to the smart contract---no raw surveillance data leaves the client unencrypted, and no raw data reaches the blockchain. This design accommodates lightweight client devices: reporters require only a browser with MetaMask and connectivity sufficient to reach an IPFS gateway and an Ethereum node endpoint.

\textbf{Tier~2: Edge/Fog Preprocessing.} The Next.js API middleware layer implements preliminary validation, schema conformance checks, and local buffering. When a reporter submits an event, the middleware validates field completeness and data-type conformance before forwarding the encrypted payload for IPFS upload. During connectivity interruptions, submissions are buffered in a local queue and retransmitted upon reconnection. This preprocessing reduces invalid submissions reaching the blockchain layer and provides near-source error feedback to reporters. In the current prototype, edge/fog logic executes within the Next.js server; production deployment would distribute these functions to dedicated edge nodes at health facility sites.

\textbf{Tier~3: District Coordination.} District and regional coordinators access a supervisory interface that aggregates validated submissions from subordinate reporters. The coordination module supports batch review: coordinators inspect submission metadata (timestamps, reporter identifiers, CIDs, and completeness status) without accessing decrypted case data unless explicitly authorized. Upon batch approval, the coordinator triggers a single anchoring transaction that registers multiple CIDs and their associated metadata on-chain, reducing per-report transaction costs. Rework requests and supervisory feedback are recorded as on-chain provenance events linked to the originating report CID.

\textbf{Tier~4: Blockchain Coordination Layer.} The smart contract layer is described in Section~\ref{subsec:contracts}. On-chain state is strictly limited to: cryptographic hashes, IPFS CID pointers, access-control lists, provenance and consent records, completeness verification state, alert events, and audit trail entries. Raw surveillance payloads are never stored on-chain.

\textbf{Tier~5: Off-Chain Encrypted Storage.} Encrypted surveillance payloads are stored on IPFS using content addressing. Each payload is encrypted at the originating client (Tier~1) before upload; the resulting CID serves as both a storage pointer and an integrity anchor. The smart contract stores only the CID and a keccak256 hash of the plaintext for independent integrity verification. Authorized retrievers obtain the CID from the smart contract, fetch the encrypted blob from IPFS, and decrypt locally using keys distributed through the permission management system. This separation ensures that neither the blockchain nor the IPFS layer alone exposes identifiable health data.

\subsection{Smart Contract Architecture}\label{subsec:contracts}

The prototype implements four interdependent smart contracts that realize the selective anchoring principle described in Section~\ref{subsec:dataflow}. Figure~\ref{fig:contracts} illustrates the contract interaction topology.

\begin{figure}[!ht]
\centering
\fbox{\parbox{0.8\columnwidth}{\centering\vspace{2em}[Smart Contract Interaction Topology]\\Figure placeholder\vspace{2em}}}
\caption{Smart contract interaction topology implementing selective anchoring, access control, and completeness verification.}\label{fig:contracts}
\end{figure}

% ===== Mermaid Source =====
% Figure: Smart Contract Interaction Topology
% Scientific purpose: Shows how the four smart contracts interact to implement
% selective anchoring, permission management, provenance tracking, and completeness
% verification without storing raw surveillance data on-chain.
%
% graph LR
%     subgraph Contracts["Smart Contract Layer (Tier 4)"]
%         RF["ReportFactory<br/>· createReport()<br/>· getReports()<br/>· anchorBatch()"]
%         RC["ReportCoordinator<br/>· anchorHash()<br/>· recordProvenance()<br/>· requestRework()<br/>· markReviewed()"]
%         ACL["AccessControl<br/>· grantAccess()<br/>· revokeAccess()<br/>· checkAccess()<br/>· setExpiry()"]
%         CV["CompletenessVerifier<br/>· registerExpected()<br/>· recordSubmission()<br/>· getCompleteness()"]
%     end
%
%     subgraph OnChain["On-Chain State"]
%         H["CID + Hash Registry"]
%         P["Permission & Consent State"]
%         PR["Provenance Trail"]
%         CS["Completeness Ratios"]
%         AE["Alert & Audit Events"]
%     end
%
%     RF -->|"deploy instance"| RC
%     RF -->|"register"| CV
%     RC --> H
%     RC --> PR
%     RC --> AE
%     ACL --> P
%     CV --> CS
%     CV -->|"gap alert"| AE
%
%     style Contracts fill:#e8e8e8,stroke:#333
%     style OnChain fill:#d0d0d0,stroke:#333
% =========================

\textbf{ReportFactory.} The factory contract serves as the system entry point. It deploys and manages \texttt{ReportCoordinator} instances for each reporting unit, maintains a registry of all deployed coordinator addresses, and supports batch anchoring of multiple submissions in a single transaction. The factory pattern ensures that each coordinator retains an unmodified copy of the contract source, preventing post-deployment tampering.

\textbf{ReportCoordinator.} Each coordinator instance manages the lifecycle of surveillance reports for a specific reporting unit. It stores only selective on-chain artifacts: the IPFS CID, a keccak256 hash of the original plaintext for integrity verification, reporter address, timestamp, and tier-level metadata. The coordinator supports provenance operations including supervisory review confirmation, rework requests with descriptive annotations, and hierarchical feedback---all recorded as immutable on-chain events. Listing~\ref{lst:anchor} illustrates the selective anchoring interface.

\begin{figure}[!ht]
\begin{verbatim}
// SPDX-License-Identifier: MIT
pragma solidity ^0.8.0;

contract ReportCoordinator {
    struct Anchor {
        string  ipfsCid;
        bytes32 contentHash;
        address reporter;
        uint256 timestamp;
        uint8   tierLevel;
        bool    reviewed;
    }

    mapping(uint256 => Anchor) public anchors;
    uint256 public anchorCount;

    event ReportAnchored(
        uint256 indexed id,
        string  ipfsCid,
        bytes32 contentHash,
        address reporter,
        uint8   tierLevel
    );
    event ReportReviewed(
        uint256 indexed id,
        address reviewer
    );
    event ReworkRequested(
        uint256 indexed id,
        address supervisor,
        string  annotation
    );

    function anchorReport(
        string  calldata _cid,
        bytes32 _hash,
        uint8   _tier
    ) external returns (uint256) {
        uint256 id = anchorCount++;
        anchors[id] = Anchor(
            _cid, _hash, msg.sender,
            block.timestamp, _tier, false
        );
        emit ReportAnchored(
            id, _cid, _hash, msg.sender, _tier
        );
        return id;
    }

    function markReviewed(uint256 _id)
        external
    {
        anchors[_id].reviewed = true;
        emit ReportReviewed(_id, msg.sender);
    }

    function requestRework(
        uint256 _id,
        string calldata _note
    ) external {
        emit ReworkRequested(
            _id, msg.sender, _note
        );
    }
}
\end{verbatim}
\caption*{Listing~1: Selective anchoring and provenance interface (\texttt{ReportCoordinator}).}\label{lst:anchor}
\end{figure}

\textbf{AccessControl.} The access-control contract implements the dynamic permission model described in Section~\ref{subsec:privacy}. Permissions are scoped by tier level, reporter--supervisor relationship, and data category. The contract supports three permission operations: \texttt{grantAccess} assigns read permission to a specific address for a given report identifier; \texttt{revokeAccess} removes a previously granted permission; and \texttt{setExpiry} attaches a time-bound deadline after which the permission automatically becomes invalid. Permission state changes are emitted as on-chain events for auditability. Listing~\ref{lst:acl} illustrates the access-control interface.

\begin{figure}[!ht]
\begin{verbatim}
contract AccessControl {
    struct Permission {
        bool    granted;
        uint256 expiry;   // 0 = no expiry
        uint8   tierScope;
    }

    // reportId => grantee => Permission
    mapping(uint256 => mapping(
        address => Permission
    )) public permissions;

    event AccessGranted(
        uint256 indexed reportId,
        address grantee,
        uint256 expiry
    );
    event AccessRevoked(
        uint256 indexed reportId,
        address grantee
    );

    function grantAccess(
        uint256 _reportId,
        address _grantee,
        uint256 _expiry,
        uint8   _tierScope
    ) external {
        permissions[_reportId][_grantee] =
            Permission(true, _expiry, _tierScope);
        emit AccessGranted(
            _reportId, _grantee, _expiry
        );
    }

    function revokeAccess(
        uint256 _reportId,
        address _grantee
    ) external {
        permissions[_reportId][_grantee]
            .granted = false;
        emit AccessRevoked(_reportId, _grantee);
    }

    function checkAccess(
        uint256 _reportId,
        address _grantee
    ) external view returns (bool) {
        Permission memory p =
            permissions[_reportId][_grantee];
        if (!p.granted) return false;
        if (p.expiry != 0 &&
            block.timestamp > p.expiry)
            return false;
        return true;
    }
}
\end{verbatim}
\caption*{Listing~2: Dynamic access control with grant, revoke, and time-bound expiry (\texttt{AccessControl}).}\label{lst:acl}
\end{figure}

\textbf{CompletenessVerifier.} The completeness contract implements the event-driven verification mechanism described in Section~\ref{subsec:completeness}. For each reporting period, expected reporting units are registered. As anchoring events arrive, the contract increments a submission counter. The completeness ratio---received submissions divided by expected units---is queryable on-chain without accessing underlying payload data. When the ratio falls below a configurable threshold at period close, the contract emits a gap alert event. Listing~\ref{lst:completeness} illustrates the completeness verification interface.

\begin{figure}[!ht]
\begin{verbatim}
contract CompletenessVerifier {
    struct Period {
        uint256 expectedCount;
        uint256 receivedCount;
        bool    closed;
    }

    // periodId => Period
    mapping(uint256 => Period)
        public periods;
    // periodId => reporter => submitted
    mapping(uint256 => mapping(
        address => bool
    )) public submitted;

    event SubmissionRecorded(
        uint256 indexed periodId,
        address reporter
    );
    event GapAlert(
        uint256 indexed periodId,
        uint256 received,
        uint256 expected
    );

    function registerPeriod(
        uint256 _periodId,
        uint256 _expectedCount
    ) external {
        periods[_periodId] = Period(
            _expectedCount, 0, false
        );
    }

    function recordSubmission(
        uint256 _periodId
    ) external {
        require(
            !submitted[_periodId][msg.sender],
            "Already submitted"
        );
        submitted[_periodId][msg.sender] = true;
        periods[_periodId].receivedCount++;
        emit SubmissionRecorded(
            _periodId, msg.sender
        );
    }

    function getCompleteness(
        uint256 _periodId
    ) external view returns (
        uint256 received, uint256 expected
    ) {
        Period memory p = periods[_periodId];
        return (p.receivedCount, p.expectedCount);
    }
}
\end{verbatim}
\caption*{Listing~3: Completeness verification with gap alerting (\texttt{CompletenessVerifier}).}\label{lst:completeness}
\end{figure}

\subsection{IPFS Integration and Off-Chain Storage Workflow}\label{subsec:ipfs_impl}

The off-chain storage workflow implements the Tier~5 data layer. Figure~\ref{fig:ipfs_workflow} illustrates the IPFS anchoring lifecycle.

\begin{figure}[!ht]
\centering
\fbox{\parbox{0.8\columnwidth}{\centering\vspace{2em}[IPFS Anchoring Workflow]\\Figure placeholder\vspace{2em}}}
\caption{IPFS anchoring workflow: encryption, upload, CID return, on-chain anchoring, and authorized retrieval.}\label{fig:ipfs_workflow}
\end{figure}

% ===== Mermaid Source =====
% Figure: IPFS Anchoring Workflow
% Scientific purpose: Shows the complete lifecycle of a surveillance payload from
% client-side encryption through IPFS upload to on-chain CID anchoring and
% authorized retrieval with decryption.
%
% sequenceDiagram
%     participant Client as Reporter Client
%     participant Crypto as AES-256-GCM
%     participant IPFS as IPFS Gateway
%     participant SC as ReportCoordinator
%     participant ACL as AccessControl
%     participant Retriever as Authorized Retriever
%
%     Client->>Crypto: Encrypt payload
%     Crypto-->>Client: ciphertext + IV + tag
%     Client->>IPFS: Upload encrypted blob
%     IPFS-->>Client: Return CID
%     Client->>Client: Compute keccak256(plaintext)
%     Client->>SC: anchorReport(CID, hash, tier)
%     SC-->>Client: Confirm anchor ID
%     Note over SC: On-chain: CID + hash + metadata
%     Note over IPFS: Off-chain: encrypted blob only
%     Retriever->>ACL: checkAccess(reportId)
%     ACL-->>Retriever: Authorized
%     Retriever->>SC: getAnchor(reportId)
%     SC-->>Retriever: CID + hash
%     Retriever->>IPFS: Fetch by CID
%     IPFS-->>Retriever: Encrypted blob
%     Retriever->>Crypto: Decrypt with authorized key
%     Crypto-->>Retriever: Plaintext payload
%     Retriever->>Retriever: Verify keccak256 == on-chain hash
% =========================

The upload workflow proceeds as follows: (1)~the reporter client encrypts the surveillance payload using AES-256-GCM with a per-report symmetric key; (2)~the encrypted ciphertext, initialization vector, and authentication tag are uploaded to IPFS via \texttt{ipfs-http-client}; (3)~IPFS returns a content identifier (CID); (4)~the client computes a keccak256 hash of the original plaintext; (5)~the client calls \texttt{anchorReport(CID, hash, tier)} on the \texttt{ReportCoordinator} contract, committing only the CID, hash, and tier-level metadata on-chain. Listing~\ref{lst:ipfs} illustrates the client-side upload workflow.

\begin{figure}[!ht]
\begin{verbatim}
const { create } = require('ipfs-http-client');
const crypto = require('crypto');
const ipfs = create({ url: IPFS_GATEWAY });

async function uploadAndAnchor(
    payload, contract, account
) {
    // 1. Encrypt payload
    const key = crypto.randomBytes(32);
    const iv  = crypto.randomBytes(12);
    const cipher = crypto.createCipheriv(
        'aes-256-gcm', key, iv
    );
    let enc = cipher.update(
        JSON.stringify(payload), 'utf8', 'hex'
    );
    enc += cipher.final('hex');
    const tag = cipher.getAuthTag().toString('hex');
    const blob = JSON.stringify({ enc, iv:
        iv.toString('hex'), tag });

    // 2. Upload to IPFS
    const { cid } = await ipfs.add(
        Buffer.from(blob)
    );

    // 3. Compute plaintext hash
    const hash = web3.utils.soliditySha3(
        JSON.stringify(payload)
    );

    // 4. Anchor on-chain
    await contract.methods.anchorReport(
        cid.toString(), hash, tierLevel
    ).send({ from: account });

    return { cid: cid.toString(), key, iv };
}
\end{verbatim}
\caption*{Listing~4: IPFS upload and selective anchoring workflow.}\label{lst:ipfs}
\end{figure}

Authorized retrieval reverses the process: the retriever queries the \texttt{AccessControl} contract to verify permission status, obtains the CID from the \texttt{ReportCoordinator}, fetches the encrypted blob from IPFS, decrypts using the authorized key, and verifies plaintext integrity against the on-chain hash. This workflow ensures that blockchain is used exclusively for coordination artifacts while bulk health data remains encrypted off-chain.

\subsection{Functional Operations}\label{subsec:operations}

The prototype supports the following surveillance operations, each mapped to the architectural workflow defined in Section~\ref{sec:model}:

\textbf{Case report creation (Tiers~1$\rightarrow$5$\rightarrow$4).} A community reporter encrypts case data at the client, uploads the encrypted payload to IPFS, and submits an anchoring transaction through MetaMask to the \texttt{ReportFactory} contract. The factory deploys a \texttt{ReportCoordinator} instance (or routes to an existing one) and records the CID, content hash, and reporter metadata on-chain. No raw case data is transmitted to the blockchain (Figures~\ref{fig:create1}--\ref{fig:create3}).

\begin{figure}[!ht]
\centering
\fbox{\parbox{0.8\columnwidth}{\centering\vspace{2em}[Case Report Creation]\\Figure placeholder\vspace{2em}}}
\caption{Case report creation: client-side encryption, IPFS upload, and selective anchoring.}\label{fig:create1}
\end{figure}

\begin{figure}[!ht]
\centering
\fbox{\parbox{0.8\columnwidth}{\centering\vspace{2em}[MetaMask Transaction]\\Figure placeholder\vspace{2em}}}
\caption{MetaMask transaction signing for anchoring transaction.}\label{fig:create2}
\end{figure}

\begin{figure}[!ht]
\centering
\fbox{\parbox{0.8\columnwidth}{\centering\vspace{2em}[Anchor Confirmation]\\Figure placeholder\vspace{2em}}}
\caption{On-chain anchor confirmation with CID reference.}\label{fig:create3}
\end{figure}

\textbf{Report listing and metadata retrieval (Tier~4).} The application queries the deployed \texttt{ReportFactory} to retrieve all registered coordinator addresses, then renders submission metadata (timestamps, tier levels, completeness status, CID references) through the React front-end. Individual anchor details are retrieved through the coordinator's view functions. Actual report content requires authorized retrieval from IPFS (Figures~\ref{fig:list},~\ref{fig:detail}).

\begin{figure}[!ht]
\centering
\fbox{\parbox{0.8\columnwidth}{\centering\vspace{2em}[Report Metadata List]\\Figure placeholder\vspace{2em}}}
\caption{Report metadata listing interface showing anchor records without raw data.}\label{fig:list}
\end{figure}

\begin{figure}[!ht]
\centering
\fbox{\parbox{0.8\columnwidth}{\centering\vspace{2em}[Report Detail with CID]\\Figure placeholder\vspace{2em}}}
\caption{Report detail display with CID reference, integrity hash, and permission status.}\label{fig:detail}
\end{figure}

\textbf{Supervisory review and provenance tracking (Tiers~3$\rightarrow$4).} District-level supervisors review submission metadata and, upon authorization via the \texttt{AccessControl} contract, retrieve and decrypt case data from IPFS. Review actions (\texttt{markReviewed}) and rework requests (\texttt{requestRework} with descriptive annotations) are recorded as immutable on-chain provenance events, creating a complete supervisory audit trail (Figures~\ref{fig:review},~\ref{fig:feedback}).

\begin{figure}[!ht]
\centering
\fbox{\parbox{0.8\columnwidth}{\centering\vspace{2em}[Supervisory Review]\\Figure placeholder\vspace{2em}}}
\caption{Supervisory review interface with provenance tracking.}\label{fig:review}
\end{figure}

\begin{figure}[!ht]
\centering
\fbox{\parbox{0.8\columnwidth}{\centering\vspace{2em}[Rework Request]\\Figure placeholder\vspace{2em}}}
\caption{Rework request with on-chain annotation and audit event.}\label{fig:feedback}
\end{figure}

\textbf{Completeness verification (Tiers~3$\rightarrow$4).} The \texttt{CompletenessVerifier} contract tracks expected versus received anchoring events for each reporting period and catchment area. The completeness ratio is computed on-chain from submission-anchor metadata, not from raw payload data. Gap alerts are emitted when the ratio falls below the configured threshold at period close (Section~\ref{subsec:completeness}).

\textbf{Surveillance data analysis (Tier~3).} Epidemiological analysis---temporal, spatial, and demographic---executes at the district coordination layer over decrypted off-chain data retrieved through permission-controlled access. Only summary indicators and alert triggers are anchored on-chain. The analysis module presents aggregated results through a dashboard interface (Figure~\ref{fig:analysis}).

\begin{figure}[!ht]
\centering
\fbox{\parbox{0.8\columnwidth}{\centering\vspace{2em}[Analysis Dashboard]\\Figure placeholder\vspace{2em}}}
\caption{Surveillance data analysis dashboard presenting temporal, spatial, and demographic indicators.}\label{fig:analysis}
\end{figure}

\textbf{Laboratory result reporting (Tiers~1$\rightarrow$5$\rightarrow$4).} Laboratory technicians encrypt pathogen identification results at the client, upload to IPFS, and submit an anchoring transaction. The workflow follows the same encrypt--upload--anchor pattern as case reports, ensuring consistent off-chain storage and on-chain coordination for all surveillance data types (Figure~\ref{fig:lab}).

\begin{figure}[!ht]
\centering
\fbox{\parbox{0.8\columnwidth}{\centering\vspace{2em}[Laboratory Reporting]\\Figure placeholder\vspace{2em}}}
\caption{Laboratory result reporting with encrypted off-chain storage and selective anchoring.}\label{fig:lab}
\end{figure}

\subsection{Transaction, Data, and Control Flow Summary}\label{subsec:flow_summary}

Table~\ref{tab:flow} summarizes the flow design across the prototype:

\begin{table}[!ht]
\caption{Transaction, data, and control flow across architectural tiers.}\label{tab:flow}
\centering
\small
\begin{tabular}{p{2.2cm}p{2.5cm}p{2.8cm}}
\toprule
\textbf{Flow Type} & \textbf{Blockchain Used} & \textbf{Blockchain Not Used} \\
\midrule
Data flow & CID pointers, content hashes, permission state & Raw surveillance payloads, case-level records, laboratory data \\
\midrule
Transaction flow & Anchoring txs, ACL updates, provenance events, completeness records, alert events & Payload encryption, IPFS upload, data analysis computation \\
\midrule
Control flow & Permission grant/revoke/expiry, report review status, rework requests & Validation preprocessing, batch aggregation, local buffering \\
\bottomrule
\end{tabular}
\end{table}

\subsection{Deployment Configuration}\label{subsec:deployment}

The prototype was deployed on a local Ethereum test network (Ganache~CLI) for development and functional verification, with contract deployment managed by a custom migration script. The deployment script instantiates the \texttt{ReportFactory} contract, which in turn deploys initial \texttt{ReportCoordinator}, \texttt{AccessControl}, and \texttt{CompletenessVerifier} instances. Listing~\ref{lst:deploy} illustrates the deployment sequence. IPFS storage was provided by a local IPFS daemon during development and testing. The Infura API provides remote Ethereum and IPFS gateway access for extended testing without requiring local full-node operation. Performance evaluation (Section~\ref{sec:evaluation}) was conducted on the Rinkeby proof-of-authority test network to measure block generation timing under realistic consensus conditions.

\begin{figure}[!ht]
\begin{verbatim}
const factory = await new web3.eth.Contract(
    ReportFactoryABI
).deploy({
    data: ReportFactoryBytecode
}).send({
    from: accounts[0],
    gas: '5000000'
});

// Factory deploys sub-contracts
const coordAddr =
    await factory.methods
    .deployCoordinator(woredaId)
    .send({ from: accounts[0] });

const aclAddr =
    await factory.methods
    .deployAccessControl()
    .send({ from: accounts[0] });

const compAddr =
    await factory.methods
    .deployCompletenessVerifier()
    .send({ from: accounts[0] });
\end{verbatim}
\caption*{Listing~5: Contract deployment sequence.}\label{lst:deploy}
\end{figure}

% ===== Mermaid Source =====
% Figure: Deployment Workflow
% Scientific purpose: Shows the one-time deployment sequence from factory
% contract to sub-contract instantiation, mapped to the Tier 4 coordination layer.
%
% flowchart TD
%     DEP["Deployer Account"]
%     RF["Deploy ReportFactory"]
%     RC["Deploy ReportCoordinator(s)"]
%     ACL["Deploy AccessControl"]
%     CV["Deploy CompletenessVerifier"]
%     REG["Contract Registry"]
%
%     DEP -->|"deploy tx"| RF
%     RF -->|"factory.deployCoordinator()"| RC
%     RF -->|"factory.deployAccessControl()"| ACL
%     RF -->|"factory.deployCompletenessVerifier()"| CV
%     RC --> REG
%     ACL --> REG
%     CV --> REG
%
%     style RF fill:#e8e8e8,stroke:#333
%     style REG fill:#d0d0d0,stroke:#333
% =========================

%% ===== EVALUATION =====
\section{Evaluation}\label{sec:evaluation}

The proposed model and its prototype implementation were evaluated across three dimensions: performance, usability, and security. This multi-dimensional approach assesses operational feasibility (performance), user acceptance (usability), and resilience to adversarial threats (security).

\subsection{Performance Evaluation}\label{subsec:performance}

Performance was assessed through block generation time measurement on the Ethereum Rinkeby test network. Ten consecutive experiments were conducted, each measuring the time to generate 12 sequential blocks. The measurement procedure operates as follows: (1) the user node queries the current block number $N$ and records the timestamp $t_N$; (2) the node repeatedly queries the current block number $n$ until 12 increments are observed; (3) at each increment, the time span between consecutive blocks is calculated as $t_n - t_N$; and (4) the block number and timestamp are updated for the next measurement cycle.

Table~\ref{tab:performance} presents representative results. Block generation times ranged from 12 to 14 seconds per block, with one outlier of 114 seconds observed in Experiment 5, Block 1 (attributable to network congestion on the shared test network). Excluding this outlier, the total time to produce 12 blocks was consistently approximately 180 seconds. Figure~\ref{fig:performance} presents the corresponding bar plot visualization.

\begin{table}[!ht]
\caption{Block generation time (seconds) -- representative experiments.}\label{tab:performance}
\begin{tabular*}{\columnwidth}{@{\extracolsep\fill}lccc@{}}
\toprule
Block & Exp. 1 & Exp. 5 & Exp. 10 \\
\midrule
1 & 14 & 114 & 12 \\
2 & 12 & 14 & 14 \\
3 & 14 & 12 & 12 \\
4 & 12 & 14 & 14 \\
5 & 14 & 12 & 12 \\
6 & 12 & 14 & 14 \\
\bottomrule
\end{tabular*}
\end{table}

\begin{figure}[!ht]
\centering
\fbox{\parbox{0.8\columnwidth}{\centering\vspace{2em}[Performance Bar Plot]\\Figure placeholder\vspace{2em}}}
\caption{Block generation time across 10 experiments.}\label{fig:performance}
\end{figure}

The narrow range of 12--14 seconds per block demonstrates stable and predictable transaction confirmation times under the Proof of Authority (PoA) consensus algorithm employed by the Rinkeby network. PoA eliminates the computational resource expenditure associated with Proof of Work, yielding consistent block generation without the variability inherent in mining-based consensus.

\subsection{Usability Evaluation}\label{subsec:usability}

Usability was assessed using the Software Usability Measurement Inventory (SUMI), an internationally standardized psychometric instrument validated for software usability measurement. An academic license was obtained for this study. SUMI evaluates six scales: efficiency, affect, helpfulness, control, learnability, and global usability. The instrument's reliability and discriminant validity have been established through extensive factor analysis across diverse software systems.\cite{bevan1995measuring}

\textbf{Participants.} Twenty-six evaluators were recruited from the Amhara Public Health Institute: 13 field epidemiologists, 7 public health workers, and 6 laboratory technicians. Participants received a brief introduction to the research topic and a general description of the prototype, then interacted with the system for 15--20 minutes before completing the SUMI questionnaire.

\textbf{Results.} SUMI scales are standardized to a population mean of 50 with a standard deviation of 10 (50/10 distribution), on a scale ranging from 10 to 73. Figure~\ref{fig:sumi_summary} presents the statistical summary. The prototype achieved scores above the population mean on all six scales, with a maximum score of 68.96 and a minimum score of 60.46. Figure~\ref{fig:sumi_ci} presents the mean scores with 95\% confidence intervals (CIs), demonstrating that all scale CIs fall entirely above the 50-point population mean---indicating statistically significant above-average usability at the 95\% confidence level. Figure~\ref{fig:sumi_box} presents median boxplots with quartile distributions and whisker ranges.

\begin{figure}[!ht]
\centering
\fbox{\parbox{0.8\columnwidth}{\centering\vspace{2em}[SUMI Statistical Summary]\\Figure placeholder\vspace{2em}}}
\caption{SUMI evaluation statistical summary.}\label{fig:sumi_summary}
\end{figure}

\begin{figure}[!ht]
\centering
\fbox{\parbox{0.8\columnwidth}{\centering\vspace{2em}[SUMI Mean Scores with 95\% CIs]\\Figure placeholder\vspace{2em}}}
\caption{SUMI mean scores with 95\% confidence intervals.}\label{fig:sumi_ci}
\end{figure}

\begin{figure}[!ht]
\centering
\fbox{\parbox{0.8\columnwidth}{\centering\vspace{2em}[SUMI Boxplots]\\Figure placeholder\vspace{2em}}}
\caption{SUMI median boxplots with quartile distributions.}\label{fig:sumi_box}
\end{figure}

These results indicate that domain experts---epidemiologists, public health workers, and laboratory technicians---rated the prototype favorably across all usability dimensions, supporting the practical viability of blockchain-based surveillance interfaces for the target user population.

\subsection{Security Evaluation}\label{subsec:security}

Security was assessed through threat modeling and risk scenario analysis. The threat model defines the following components:

\textbf{Assets.} The primary asset is infectious disease surveillance data---high-volume, multi-source data collected and analyzed across the public health hierarchy.

\textbf{Threats.} Threats encompass scenarios in which surveillance data is exposed to unauthorized parties, altered without authorization, or prevented from reaching designated recipients.

\textbf{Attack scenarios and mitigations.} Three attack types were analyzed:

\textit{Man-in-the-Middle (MITM) Attack.} An adversary intercepts communication between model users to modify or corrupt surveillance data in transit. The proposed model mitigates MITM attacks through decentralized identity management and cryptographic hashing of all transmitted data. Each user relationship employs multi-factor authentication combining blockchain addresses and cryptographic keys distributed across participating nodes. If a user address is compromised, any fraudulent transaction will fail hash verification against the blockchain record, and the distributed architecture eliminates single-point-of-failure exploitation.

\textit{Denial of Service (DoS) Attack.} An adversary overwhelms the system with false requests to disrupt surveillance operations. The model mitigates DoS attacks through two mechanisms: (1) the Ethereum network's tokenized transaction costs impose economic barriers on flood attacks, as each transaction requires expenditure of limited-supply tokens, making sustained flooding economically infeasible against the total consortium computing capacity; and (2) a data volume agreement protocol detects anomalous transmission volumes and triggers address replacement and stream reset upon agreement violation.

\textit{Traffic Analysis Attack.} An adversary on a remote node sniffs surveillance data traffic to gain unauthorized access to report contents. The model mitigates this attack through public-key encryption of all surveillance data stored in smart contracts---data encrypted with the public key $k_{rp}$ can only be decrypted with the corresponding private key $k_{rs}$, rendering intercepted traffic unintelligible to unauthorized nodes.

\textbf{Summary.} The threat model analysis demonstrates that the proposed blockchain-based architecture provides defense-in-depth against the three primary attack vectors relevant to disease surveillance systems, with mitigation mechanisms derived directly from blockchain's inherent properties of decentralization, immutability, and cryptographic security.

%% ===== DISCUSSION =====
\section{Discussion}\label{sec:discussion}

The evaluation results validate the proposed blockchain-based DEWS model across three complementary dimensions, each addressing a specific subset of the limitations identified in Section~\ref{sec:introduction}.

\textbf{Performance implications.} The observed block generation time of 12--14 seconds per block on the Rinkeby PoA network demonstrates that blockchain-based surveillance transactions can achieve near-real-time confirmation---well within the operational requirements of disease reporting workflows, where immediate-report conditions mandate notification within 30 minutes and weekly conditions require same-day submission. This performance compares favorably with the transaction delay limitations reported by Griggs et al.,\cite{griggs2018healthcare} whose IoT-focused blockchain system identified latency as a constraint on real-time effectiveness. The PoA consensus mechanism, by eliminating mining competition, provides predictable throughput suitable for institutional health information systems where validator identity is known and trusted.

\textbf{Usability implications.} The SUMI evaluation demonstrates statistically significant above-average usability across all six scales (minimum score 60.46, maximum 68.96, population mean 50), indicating that domain professionals---epidemiologists, public health workers, and laboratory technicians---can effectively interact with blockchain-based surveillance interfaces. This finding is particularly significant given that blockchain technology introduces unfamiliar interaction patterns (wallet-based authentication, transaction confirmation) to users accustomed to conventional web applications. The results suggest that appropriate interface design can abstract blockchain complexity sufficiently for non-technical health professionals.

\textbf{Security implications.} The threat model analysis demonstrates that the proposed architecture provides inherent defense against the three primary attack vectors (MITM, DoS, traffic analysis) through blockchain's native properties---decentralized identity management, tokenized transaction costs, and cryptographic data protection. Unlike the centralized architectures of Omar et al.\cite{omar2019privacy} and Quaini et al.,\cite{quaini2018model} which rely on perimeter security, the proposed model distributes trust across all participating nodes, eliminating single points of compromise. The multi-factor authentication combining blockchain addresses with distributed cryptographic keys provides stronger identity assurance than the certificate-based approaches vulnerable to MITM attacks in centralized systems.

\textbf{Positioning against related work.} The proposed model addresses a gap not covered by prior blockchain-in-healthcare studies. Omar et al.\cite{omar2019privacy} and Zhang et al.\cite{zhang2018fhirchain} targeted patient records and clinical data sharing respectively but did not address multi-tier hierarchical reporting workflows. Griggs et al.\cite{griggs2018healthcare} focused narrowly on medical IoT security without implementing surveillance business logic. The ML-based surveillance approaches of Bollig et al.,\cite{bollig2020machine} Harvey et al.,\cite{harvey2021predicting} and Rubio-Solis et al.\cite{rubiosolis2019zika} provide predictive analytics but cannot guarantee data integrity, decentralization, or auditability---properties that the proposed blockchain model inherently provides. The present work is, to the best of available knowledge, the first to combine blockchain smart contract architecture with the complete DEWS workflow including hierarchical reporting, laboratory integration, automated completeness verification, and multi-dimensional epidemiological analysis.

\textbf{Limitations.} Several limitations constrain the generalizability of the findings. First, the prototype was deployed on the Rinkeby test network rather than the Ethereum mainnet, and test network conditions may not fully replicate production performance characteristics. Second, the usability evaluation sample ($n=26$) was drawn from a single regional institute; broader multi-site evaluation would strengthen external validity. Third, the blockchain suitability survey ($n=10$) employed purposive sampling from one institution, limiting the representativeness of stakeholder perspectives. Fourth, the security evaluation relies on theoretical threat modeling rather than penetration testing against the deployed prototype. Fifth, scalability under high transaction volumes (multiple regions reporting simultaneously) was not empirically assessed.

%% ===== CONCLUSION =====
\section{Conclusion and Future Work}\label{sec:conclusion}

This paper presented a blockchain-based architectural model for disease early warning and surveillance, implemented as a functional prototype on the Ethereum platform. The model integrates four core surveillance processes---hierarchical disease case reporting, laboratory result reporting, automated report completeness verification, and multi-dimensional epidemiological data analysis---within a decentralized, immutable, and traceable smart contract architecture.

The three-dimensional evaluation validated the model against operational requirements. Performance evaluation demonstrated stable block generation times of 12--14 seconds under Proof of Authority consensus, satisfying the temporal requirements of surveillance reporting workflows. Usability evaluation via the standardized SUMI instrument ($n=26$ domain professionals) yielded scores of 60.46--68.96 across all six scales, significantly exceeding the population mean and confirming practical usability for epidemiologists, public health workers, and laboratory technicians. Security evaluation through threat modeling demonstrated inherent resistance to man-in-the-middle, denial-of-service, and traffic analysis attacks through blockchain's native cryptographic and decentralization properties.

The proposed model addresses the identified limitations of existing disease surveillance systems---data centralization, security vulnerabilities, traceability gaps, manual completeness verification, and inter-organizational interoperability barriers---through the systematic application of distributed ledger technology to the public health domain.

Future research directions include:
\begin{itemize}
\item Extension of the model to encompass public health emergency response and recovery processes downstream of surveillance.
\item Implementation of incidence proportion and prevalence rate computation functions within the smart contract architecture.
\item Upgrade of the prototype from test network deployment to a production-grade operational system.
\item Evaluation of alternative blockchain platforms beyond the three compared in this study.
\item Integration of machine learning and artificial intelligence modules with the blockchain architecture for predictive surveillance capabilities.
\item Development of a cross-platform mobile application to improve field accessibility.
\item Comparative evaluation against similar blockchain-based health information systems.
\end{itemize}

%% ===== BACK MATTER =====

\section*{Author Contributions}
\textbf{Teklhiwot Mekonen:} Conceptualization; Methodology; Software; Validation; Formal analysis; Investigation; Data curation; Writing -- original draft; Visualization.

\textbf{Esubalew Alemneh:} Supervision; Writing -- review \& editing.

\section*{Acknowledgments}
The authors thank the staff of the Amhara Public Health Institute for their participation in the blockchain suitability assessment survey and the SUMI usability evaluation. The authors also acknowledge the Human Computer Interaction Research Group (HCIRG), University College Cork, Ireland, for provision of the SUMI academic license.

\section*{Financial Disclosure}
This research was conducted as part of a Master of Science thesis at Bahir Dar University, Bahir Dar Institute of Technology, Faculty of Computing. No specific external grant from any funding agency in the public, commercial, or not-for-profit sectors was received for this research.

\section*{Conflict of Interest}
The authors declare no conflicts of interest.

\section*{Data Availability Statement}
The smart contract source code for the prototype is available in the repository accompanying this manuscript. The SUMI questionnaire instrument is proprietary and available under academic license from the Human Computer Interaction Research Group, University College Cork, Ireland. Raw SUMI response data are not publicly available due to participant confidentiality but are available from the corresponding author upon reasonable request.

\section*{Ethics Statement}
This study involved a usability evaluation with 26 human participants (13 field epidemiologists, 7 public health workers, and 6 laboratory technicians) recruited from the Amhara Public Health Institute, Bahir Dar, Ethiopia. Participants were briefed on the research objectives and provided voluntary agreement to participate. The study involved interaction with a software prototype for 15--20 minutes followed by completion of the standardized Software Usability Measurement Inventory (SUMI) questionnaire. The study posed no more than minimal risk to participants: no personal health data were collected, no clinical interventions were performed, and all questionnaire responses were anonymous.

\bibliography{references}

\end{document}
